\subsubsection{Regulatory and Structural Evolution}

\begin{principle} \hlt{Regulation as Market Design}\\
Regulation is market design by law. Good rules support communication, performance, investor protection, competition, and access; bad rules protect incumbents or raise trading costs.
\end{principle}

\begin{taxonomy} \hlt{Market-Design Question Segmentation}
\begin{enumerate}[label=\roman*.]
\setlength{\itemsep}{0pt}
\item \hlt{consolidation}: central book versus fragmented venues;
\item \hlt{intermediation}: order-driven, quote-driven, brokered, or hybrid;
\item \hlt{access}: internalisation, preferencing, linkages, and who may trade;
\item \hlt{transparency}: order-book visibility, trade reporting, and market-data ownership;
\item \hlt{controls}: tick sizes, price limits, halts, and listing permissions.
\end{enumerate}
\end{taxonomy}

\begin{process} \hlt{Welfare Evaluation}
\begin{enumerate}[label=\roman*.]
\setlength{\itemsep}{0pt}
\item identify feasible market structures;
\item state the objective function and trade-offs;
\item estimate effects on liquidity, information, volatility, and distribution;
\item choose the rule with the best expected social welfare.
\end{enumerate}
\end{process}

\begin{taxonomy} \hlt{Policy Objective Segmentation}
\begin{enumerate}[label=\roman*.]
\setlength{\itemsep}{0pt}
\item \hlt{utilitarian traders}: protect the users whose needs make markets exist;
\item \hlt{public benefits}: promote informative prices and low-cost liquidity;
\item \hlt{profit-motivated traders}: support them only when they supply liquidity or information;
\item \hlt{predatory activity}: discourage manipulation, bluffing, and front-running.
\end{enumerate}
\end{taxonomy}

\begin{taxonomy} \hlt{Exposure-Protection Rule Segmentation}
\begin{enumerate}[label=\roman*.]
\setlength{\itemsep}{0pt}
\item \hlt{time precedence}: protects displayed orders from later traders at the same price;
\item \hlt{minimum price increment}: makes stepping ahead costly enough to preserve queue value;
\item \hlt{undisplayed orders}: let traders make firm commitments without full pre-trade exposure;
\item \hlt{audit trail}: deters broker front-running, dual-capacity abuse, and allocation manipulation;
\item \hlt{delayed block reporting}: can protect large trades while they are being completed.
\end{enumerate}
These reduce exposure costs, but each also trades off public transparency, price discovery, or queue competition.
\end{taxonomy}

\begin{marketrule} \hlt{Decimalization}\\
Decimalization reduced U.S. equity tick sizes from fractions to cents. Smaller ticks narrowed quoted spreads but reduced displayed depth and queue stability.
\end{marketrule}

\begin{marketrule} \hlt{Reg NMS}\\
Reg NMS links U.S. equity venues through protected quotes and routing obligations. It protects displayed top-of-book prices, while depth, fees, latency, and hidden liquidity still differ across venues.
\end{marketrule}

\begin{marketrule} \hlt{MiFID and MiFID II}\\
MiFID broke exchange monopolies by allowing trading outside historical exchanges; MiFID II tightened best-execution, transparency, and reporting duties. The combined effect was more competition, more venue types, and greater reliance on routing, reporting, and post-trade data quality.
\end{marketrule}

\begin{principle} \hlt{Best Execution After Fragmentation}\\
Best execution is not only the best displayed price. In fragmented markets it also depends on fees, likelihood of execution, latency, hidden liquidity, venue access, order size, and the investor's horizon.
\end{principle}

\begin{mechanism} \hlt{Fragmentation}\\
Fragmentation means the same instrument trades across multiple venues. It can lower fees and specialise execution, but splits liquidity and complicates best execution.
\end{mechanism}

\begin{taxonomy} \hlt{Levels of Fragmentation}
\begin{enumerate}[label=\roman*.]
\setlength{\itemsep}{0pt}
\item \hlt{operators}: new market operators compete with incumbents, then often merge as fixed costs bite;
\item \hlt{venues}: one operator may run lit books, dark books, periodic auctions, and internalisation channels;
\item \hlt{order books}: one venue may maintain several books with different access, transparency, or priority rules;
\item \hlt{orders}: large parent orders are split through time and across compatible venues;
\item \hlt{technology}: protocols, market data feeds, and native gateways become part of venue competition.
\end{enumerate}
\end{taxonomy}

\begin{mechanism} \hlt{Structural Evolution of Fragmentation}\\
Post-MiFID fragmentation first reduced the monopoly share of historical exchanges, then consolidated around large pan-European operators. Complexity did not disappear: liquidity was segmented into lit books, dark pools, broker crossing systems, systematic internalisers, and specialised order books.
\end{mechanism}

\begin{taxonomy} \hlt{Post-MiFID Phases}
\begin{enumerate}[label=\roman*.]
\setlength{\itemsep}{0pt}
\item \hlt{entry}: new MTFs competed for order flow, often with bank sponsorship or imported technology;
\item \hlt{convergence}: successful venues moved toward pan-European offers and common functionality;
\item \hlt{dark segmentation}: broker crossing networks, dark MTFs grew around first-hand liquidity, lower signalling;
\item \hlt{consolidation}: fixed technology, data, and connectivity costs pushed operators toward mergers.
\end{enumerate}
\end{taxonomy}

\begin{caveat} \hlt{Early MiFID Failures}\\
Turquoise and NASDAQ OMX Europe show that technology and sponsorship were not sufficient. Turquoise gained relevance only after its liquidity agreement ended and it became part of the LSE group; NASDAQ OMX's European MTF and router never reached durable scale.
\end{caveat}

\begin{caveat} \hlt{Liquidity Segmentation}\\
Venue competition separates natural investor liquidity from opportunistic market-making liquidity. This can improve matching for specific order types, but the trader now pays for aggregation through routing logic, data, connectivity, and information leakage.
\end{caveat}

\begin{taxonomy} \hlt{Technology Fragmentation}
\begin{enumerate}[label=\roman*.]
\setlength{\itemsep}{0pt}
\item \hlt{application protocols}: venue APIs and market-data formats such as ITCH, XDP, EMDI, or RWF;
\item \hlt{network protocols}: multicast UDP, TCP, replay channels, and venue-specific recovery logic;
\item \hlt{compression protocols}: FAST, ZLIB, or uncompressed feeds.
\end{enumerate}
Even when trading rules converge, access remains fragmented through protocols, feeds, and operational detail.
\end{taxonomy}

\begin{marketrule} \hlt{Transparency and Dark Trading}\\
Pre-trade transparency reveals resting interest; post-trade transparency reveals completed transactions. Dark trading reduces information leakage, so regulators often pair it with caps, waivers, or stronger reporting.
\end{marketrule}

\begin{mechanism} \hlt{Broker Crossing Networks}\\
Broker crossing networks internalise client flow before exposing it to lit markets.\\
They can protect large investors from signalling, but they also raise conflict questions when broker prop flow, market-making flow, and client flow meet in the same pool.
\end{mechanism}

\begin{marketrule} \hlt{MiFID II Dark Caps}\\
MiFID II constrained dark trading through double volume caps: a venue-level cap and an aggregate market-level cap. The intent is to preserve price formation on lit venues while still allowing limited dark execution.
\end{marketrule}

\begin{taxonomy} \hlt{MiFID Venue Statuses}\\
European transactions may occur on regulated markets, multilateral trading facilities, systematic internalisers, or OTC channels. Regulated markets and MTFs use multilateral rule books; systematic internalisers trade bilaterally against client flow; OTC trading keeps negotiation bilateral but still faces reporting obligations.
\end{taxonomy}

\begin{marketrule} \hlt{Tick Size Reforms}\\
Tick size rules balance price competition against queue discipline. Large ticks widen spreads but support displayed queues; small ticks narrow spreads but encourage queue jumping, cancellation, smaller displayed size.
\end{marketrule}

\begin{caveat} \hlt{Tick Size as Market Share Tool}\\
In fragmented markets, venues can use smaller ticks to undercut incumbent quotes and attract flow. This may improve displayed prices in the short run, but can also start a race to economically meaningless price increments.
\end{caveat}

\begin{marketrule} \hlt{Sub-Penny and Tick-Size Wars}\\
Sub-penny rules restrict displayed price improvement, but off-exchange venues can still internalise flow at tiny increments. In Europe, early post-MiFID venues used different tick tables to compete until harmonisation reduced tick-size competition.
\end{marketrule}

\begin{taxonomy} \hlt{Who Gains From Tick Changes}
\begin{enumerate}[label=\roman*.]
\setlength{\itemsep}{0pt}
\item \hlt{liquidity takers}: may gain from narrower spreads on small trades;
\item \hlt{liquidity providers}: lose queue rents when stepping ahead becomes cheap;
\item \hlt{large investors}: may face less displayed depth and higher search or slicing costs;
\item \hlt{venues}: can gain share by offering finer price grids or specialised order types.
\end{enumerate}
\end{taxonomy}

\begin{principle} \hlt{Tick Policy Design}\\
Tick tables should vary with price and liquidity, stay stable enough for adaptation, and be revised using spreads, depth, order-to-trade ratios, and standard-size execution costs.
\end{principle}

\begin{marketrule} \hlt{Circuit Breakers and Structural Controls}\\
Halts, limit moves, short-sale constraints, margin rules, and position limits do not remove volatility. They change when and how traders express demand, supply, and leverage.
\end{marketrule}

\begin{taxonomy} \hlt{Circuit Breaker Scope}
\begin{enumerate}[label=\roman*.]
\setlength{\itemsep}{0pt}
\item \hlt{local}: halt one instrument or one venue;
\item \hlt{related-instrument}: halt derivatives, ETFs, or baskets tied to the halted instrument;
\item \hlt{market-wide}: halt a broad market after index-level stress;
\item \hlt{dynamic bands}: reject or pause trades outside a price corridor;
\item \hlt{static thresholds}: trigger from a fixed reference price or previous close.
\end{enumerate}
\end{taxonomy}

\begin{principle} \hlt{Circuit Breaker Design}\\
Circuit breakers should stop disorderly liquidity demand, not ordinary price discovery. Useful triggers combine price movement with liquidity, volatility, correlation, or related-instrument stress.
\end{principle}

\begin{caveat} \hlt{Propagation Problem}\\
Halting one instrument can push pressure into substitutes. If ETFs, futures, options, or highly correlated names remain open, the halt may move price discovery elsewhere instead of calming it.
\end{caveat}

\begin{marketrule} \hlt{Post-Flash-Crash Controls}\\
After the 2010 Flash Crash, U.S. market structure moved toward single-stock pauses, limit-up/limit-down bands, consolidated audit trail reporting, stub-quote restrictions, and tighter controls on naked market access.
\end{marketrule}

\begin{caveat} \hlt{Stub Quotes}\\
Stub quotes satisfy a formal two-sided quoting duty at prices far from the market. They are not real liquidity; executions against them reveal missing protection in the trading mechanism.
\end{caveat}

\begin{principle} \hlt{Consolidated Tape as Resilience}\\
Helps fragmented markets retain a shared post-trade reference. Without it, stress or outages can leave participants unsure whether price discovery is happening in one market or several disconnected ones.
\end{principle}

\begin{taxonomy} \hlt{Regulatory Institutions}\\
Market oversight is shared across:
\begin{enumerate}[label=\roman*.]
\setlength{\itemsep}{0pt}
\item governmental regulators, such as securities and futures commissions;
\item self-regulatory organisations, such as exchanges, clearinghouses, and trader associations;
\item private standard setters, such as accounting, performance-reporting, and professional bodies;
\item international organisations, such as IOSCO and the World Federation of Exchanges.
\end{enumerate}
This division creates expertise, jurisdictional overlap, and regulatory competition.
\end{taxonomy}

\begin{caveat} \hlt{Unintended Consequences}\\
Rules change incentives. Taxes may move liquidity abroad, bans may move risk into substitutes, protected quotes may increase routing complexity, and exchange rules may protect incumbents.
\end{caveat}

\begin{mechanism} \hlt{Competition to Merger}\\
Opening trading to competition first produces innovative venues and lower explicit costs. Once practices converge, scale becomes decisive: mergers reduce fixed technology costs, expand connectivity, and preserve distinct order books under larger operating groups.
\end{mechanism}

\begin{practitionernote} \hlt{Electronification Beyond Equities}\\
Electronification does not imply one architecture. Equities move to lit books and dark pools; fixed income and derivatives retain RFQ, dealer networks, and central clearing because contracts and counterparty risks differ.
\end{practitionernote}
